% part: intuitionistic-logic
% chapter: propositions-as-types
% section: sequent-natural-deduction

\documentclass[../../../include/open-logic-section]{subfiles}

\begin{document}

\olfileid[pt]{int}{pty}{snd}

\olsection{Dedução natural em sequentes}

Escrevemos $\Gamma \Sequent !A$ quando existe uma !!{derivation} de
dedução natural com $\Gamma$ como hipóteses !!{undischarged} e $!A$
como conclusão; ou $\Sequent !A$ quando $\Gamma$ é vazio.

Escrevemos $\Gamma, !A_1, \dots, !A_n$ para $\Gamma \cup \{!A_1, \dots,
!A_n\}$ e $\Gamma, \Delta$ para $\Gamma \cup \Delta$.

Observe que, quando temos $\Gamma \Sequent !A \land !A$, isso significa
que temos !!a{derivation} com $\Gamma$ como hipóteses !!{undischarged} e
$!A \land !A$ como !!{formula} final. Aplicando $\Elim{\land}$ na parte
inferior, podemos obter !!a{derivation} com as mesmas hipóteses
!!{undischarged} e $!A$ como conclusão. Em outras palavras, se $\Gamma
\Sequent !A \land !B$, então $\Gamma \Sequent !A$.
\begin{prooftree}
  \Axiom$\Gamma \fCenter !A \land !B$
  \RightLabel{$\Elim{\land}$}
  \UnaryInf$\Gamma \fCenter !A$
  \DisplayProof\qquad\bottomAlignProof
  \Axiom$\Gamma \fCenter !A \land !B$
  \RightLabel{$\Elim{\land}$}
  \UnaryInf$\Gamma \fCenter !B$
\end{prooftree}
O rótulo $\Elim{\land}$ sugere a relação com a regra de mesmo nome na
dedução natural.

De modo semelhante, suponha que temos $\Gamma, !A \Sequent !B$, isto é,
!!a{derivation} com as hipóteses !!{undischarged} $\Gamma, !A$ e a
!!{formula} final~$!B$. Se aplicarmos a regra $\Intro{\lif}$, teremos
!!a{derivation} com $\Gamma$ como hipóteses !!{undischarged} e $!A \lif
!B$ como !!{formula} final, isto é, $\Gamma \Sequent !A \lif !B$.
Observe como isso torna mais explícita a !!{discharge} de hipóteses.
\begin{prooftree}
  \Axiom $\Gamma, !A \fCenter !B$
  \RightLabel{$\Intro{\lif}$}
  \UnaryInf $\Gamma \fCenter !A \lif !B$
\end{prooftree}

Podemos extrair conclusões das outras regras da mesma maneira, conforme
detalhado a seguir:
\begin{gather*}
  \Axiom $\Gamma \fCenter !A$
  \Axiom $\Delta \fCenter !B$
  \RightLabel{$\Intro{\land}$}
  \BinaryInf $\Gamma,\Delta \fCenter !A \land !B$
  \DisplayProof\\
  \Axiom $\Gamma \fCenter !A \land !B$
  \RightLabel{$\Elim{\land}_1$}
  \UnaryInf $\Gamma \fCenter !A$
  \DisplayProof
  \qquad
  \Axiom $\Gamma \fCenter !A \land !B$
  \RightLabel{$\Elim{\land}_2$}
  \UnaryInf $\Gamma \fCenter !B$
  \DisplayProof
  \\
  \Axiom $\Gamma \fCenter !A$
  \RightLabel{$\Intro{\lor}_1$}
  \UnaryInf $\Gamma \fCenter !A \lor !B$
  \DisplayProof\qquad
  \Axiom $\Gamma \fCenter !B$
  \RightLabel{$\Intro{\lor}_2$}
  \UnaryInf $\Gamma \fCenter !A \lor !B$
  \DisplayProof\\
  \Axiom $\Gamma \fCenter !A \lor !B$
  \Axiom $\Delta, !A \fCenter !C$
  \Axiom $\Delta', !B \fCenter !C$
  \RightLabel{$\Elim{\lor}$}
  \TrinaryInf $\Gamma, \Delta, \Delta' \fCenter !C$
  \DisplayProof
  \\
  \Axiom $\Gamma, !A \fCenter !B$
  \RightLabel{$\Intro{\lif}$}
  \UnaryInf $\Gamma \fCenter !A \lif !B$
  \DisplayProof
  \qquad
  \Axiom $\Delta \fCenter !A \lif !B$
  \Axiom $\Gamma \fCenter !A$
  \RightLabel{$\Elim{\lif}$}
  \BinaryInf $\Gamma, \Delta \fCenter !B$
  \DisplayProof
  \\
  \Axiom $\Gamma \fCenter \lfalse$
  \RightLabel{$\FalseInt$}
  \UnaryInf $\Gamma \fCenter !A$
  \DisplayProof
\end{gather*}

Toda hipótese, por si só, é !!a{derivation} de $!A$ a partir de~$!A$,
isto é, sempre temos $!A \Sequent !A$.

\begin{prooftree}
  \AxiomC{$ $}
  \UnaryInfC{$!A \fCenter !A$}
\end{prooftree}

Em conjunto, essas regras podem ser tomadas como um cálculo sobre quais
!!{derivation}s de dedução natural existem. Também podem ser tomadas
como uma variante notacional da dedução natural, na qual cada passo
registra não apenas a !!{formula} !!{derive}d, mas também as hipóteses
!!{undischarged} a partir das quais ela foi !!{derive}d.

\begin{prooftree}
  \Axiom$!A \fCenter !A$
  \UnaryInf $!A \fCenter !A \lor (!A \lif \lfalse)$
  \Axiom $!B \fCenter !B$
  \BinaryInf$!A, !B\lif \fCenter \lfalse$
  \UnaryInf$(!B \fCenter !A \lif \lfalse$
  \UnaryInf$( !B \fCenter !A \lor (!A \lif \lfalse)$
  \Axiom$(!B \fCenter !B$
  \BinaryInf$(!B \fCenter \lfalse$
  \UnaryInf$\fCenter !B \lif \lfalse$
\end{prooftree}
em que $!B$ abrevia $(!A \lor (!A \lif \lfalse))\lif \lfalse$.

\end{document}
